% Zumdahl Chapter 11 test -- CED content only (11.1-11.3).
% 20 MCQ + 2 FRQ. 50 min.
\documentclass{apchem-exam}

\apcunitword{Chapter}
\apcunit{11}
\apcunittitle{Properties of Solutions}
\apcdoctitle{Chapter Test}
\apcsubtitle{Zumdahl \S11.1--11.3 only \textbullet\ 50 minutes \textbullet\ see scope note}

\begin{document}
\apctitleblock

\begin{apcnote}[Scope]
This test covers \S11.1--11.3 only. Nothing from \S11.4--11.8 --- vapour
pressure lowering, boiling-point elevation, freezing-point depression,
osmotic pressure, or colloids --- appears. Those were marked enrichment in
the notes and are not on the AP framework.
\end{apcnote}

\examsection{I}{Multiple Choice}{20 questions \textbullet\ 30 minutes \textbullet\ 1 point each}{\nocalculator}

% ---- 11.1 composition ------------------------------------------------------
\begin{mcq}
  Molarity is defined as moles of solute per
  \choices
    \choice{kilogram of solvent}
    \correct{litre of solution}
    \choice{litre of solvent}
    \choice{kilogram of solution}
\end{mcq}
\why{Solution, not solvent --- dissolving changes the volume.}
\distractor{C}{this confusion is why volumetric flasks are filled TO the
mark}

\begin{mcq}
  Molality is defined as moles of solute per
  \choices
    \correct{kilogram of solvent}
    \choice{litre of solution}
    \choice{kilogram of solution}
    \choice{mole of solvent}
\end{mcq}
\why{Mass-based, which is what makes it temperature-independent.}

\begin{mcq}
  A solution is warmed from \SI{20}{\celsius} to \SI{50}{\celsius}. Which
  quantity changes?
  \choices
    \correct{molarity}
    \choice{molality}
    \choice{mass percent}
    \choice{mole fraction}
\end{mcq}
\why{Only molarity has a volume in its denominator, and volume expands.}

\begin{mcq}
  A solution contains \SI{0.500}{\mole} of solute in \SI{250.}{\gram} of
  water. Its molality is
  \choices
    \choice{\SI{0.500}{\mole\per\kilo\gram}}
    \correct{\SI{2.00}{\mole\per\kilo\gram}}
    \choice{\SI{1.25}{\mole\per\kilo\gram}}
    \choice{\SI{125}{\mole\per\kilo\gram}}
\end{mcq}
\why{$0.500/0.250 = 2.00$.}

\begin{mcq}
  Which concentration unit converts most directly into a partial pressure?
  \choices
    \choice{molarity}
    \choice{molality}
    \correct{mole fraction}
    \choice{mass percent}
\end{mcq}
\why{$P_A = \chi_A P_{\text{total}}$.}

\begin{mcq}
  \SI{25.0}{\gram} of solute is dissolved to give \SI{125.0}{\gram} of
  solution. The mass percent of solute is
  \choices
    \choice{25.0\%}
    \correct{20.0\%}
    \choice{16.7\%}
    \choice{80.0\%}
\end{mcq}
\why{$25.0/125.0 \times 100$ --- the denominator is the solution mass.}
\distractor{A}{used the solvent mass of 100.0 g as the denominator}

% ---- 11.2 energetics -------------------------------------------------------
\begin{mcq}
  In the three-step model of solution formation, which step is exothermic?
  \choices
    \choice{separating the solute}
    \choice{expanding the solvent}
    \correct{allowing solute and solvent to interact}
    \choice{all three}
\end{mcq}
\why{Steps 1 and 2 break attractions; only step 3 forms them.}

\begin{mcq}
  For an ionic solute dissolving in water, step 1 of the model corresponds
  to overcoming the
  \choices
    \correct{lattice energy}
    \choice{hydration energy}
    \choice{ionization energy}
    \choice{bond energy}
\end{mcq}
\why{Separating the solute means pulling the ionic lattice apart.}

\begin{mcq}
  Dissolving \ce{NH4NO3} in water makes the container feel cold. Therefore
  \choices
    \choice{$\Delta H_{\text{soln}}$ is negative}
    \correct{$\Delta H_{\text{soln}}$ is positive}
    \choice{no energy is transferred}
    \choice{the process is impossible}
\end{mcq}
\why{The system absorbed energy from the surroundings --- endothermic.}

\begin{mcq}
  A salt has $\Delta H_{\text{soln}} = \SI{-82}{\kilo\joule\per\mole}$. When
  it dissolves, the hydration energy
  \choices
    \correct{exceeds the combined cost of steps 1 and 2}
    \choice{is smaller than the lattice energy}
    \choice{is zero}
    \choice{is endothermic}
\end{mcq}
\why{A negative overall value means step 3 more than repaid steps 1 and 2.}

\begin{mcq}
  Oil does not dissolve in water primarily because
  \choices
    \choice{oil molecules are too large}
    \correct{the weak oil--water attractions cannot repay the energy of
      breaking water's hydrogen bonds}
    \choice{oil is denser than water}
    \choice{oil has no intermolecular forces}
\end{mcq}
\why{Step 2 is very costly and step 3 offers only dipole--induced dipole
attraction.}
\distractor{D}{oil molecules have dispersion forces}

\begin{mcq}
  A substance dissolves even though $\Delta H_{\text{soln}}$ is positive.
  The best explanation is that
  \choices
    \choice{the measurement was wrong}
    \correct{the increase in disorder on mixing can drive the process}
    \choice{no bonds were broken}
    \choice{the solvent was heated}
\end{mcq}
\why{Enthalpy alone does not determine whether a process occurs.}

\begin{mcq}
  \ce{NaCl} has $\Delta H_{\text{soln}} = \SI{+3.9}{\kilo\joule\per\mole}$
  despite a very large lattice energy. This small value indicates that
  \choices
    \correct{the hydration energy nearly cancels the lattice energy}
    \choice{the lattice energy is small}
    \choice{no hydration occurs}
    \choice{\ce{NaCl} is insoluble}
\end{mcq}
\why{A small net is the difference between two large opposing terms.}

% ---- 11.3 solubility factors -----------------------------------------------
\begin{mcq}
  Which is most soluble in water?
  \choices
    \choice{\ce{CCl4}}
    \choice{\ce{C6H14}}
    \correct{\ce{CH3OH}}
    \choice{\ce{I2}}
\end{mcq}
\why{Methanol hydrogen bonds with water; the others are nonpolar.}

\begin{mcq}
  Water solubility of straight-chain alcohols decreases as the chain
  lengthens because
  \choices
    \choice{the O--H group becomes less polar}
    \correct{the nonpolar hydrocarbon portion grows relative to the polar
      group}
    \choice{the molecules become ionic}
    \choice{hydrogen bonding stops entirely}
\end{mcq}
\why{One --OH cannot solvate an ever-longer hydrophobic tail.}

\begin{mcq}
  Vitamin A is stored in fatty tissue while vitamin C is excreted. This
  difference arises because vitamin A is
  \choices
    \correct{nonpolar and therefore fat-soluble}
    \choice{polar and therefore water-soluble}
    \choice{ionic}
    \choice{a larger molecule}
\end{mcq}
\why{Its structure is almost entirely C and H.}

\begin{mcq}
  Increasing the pressure above a solution has the greatest effect on the
  solubility of
  \choices
    \choice{a solid}
    \correct{a gas}
    \choice{a liquid}
    \choice{an ionic compound}
\end{mcq}
\why{Only gases are significantly compressible.}

\begin{mcq}
  As the temperature of water increases, the solubility of a typical gas
  \choices
    \choice{increases}
    \correct{decreases}
    \choice{is unchanged}
    \choice{first increases, then decreases}
\end{mcq}
\why{Which is why warm rivers hold less dissolved oxygen.}

\begin{mcq}
  An opened bottle of soda goes flat because
  \choices
    \choice{the \ce{CO2} reacts with water}
    \correct{the \ce{CO2} partial pressure above the liquid drops sharply}
    \choice{the temperature falls}
    \choice{the soda becomes saturated}
\end{mcq}
\why{Atmospheric \ce{CO2} partial pressure is roughly
$4\times10^{-4}$ atm versus about 5 atm when sealed.}

\begin{mcq}
  Which statement about the temperature dependence of solid solubility is
  correct?
  \choices
    \choice{All solids become more soluble when heated}
    \correct{Most solids become more soluble when heated, but some become
      less soluble}
    \choice{All solids become less soluble when heated}
    \choice{Temperature has no effect on solid solubility}
\end{mcq}
\why{Sodium sulfate is the standard counterexample.}
\distractor{A}{tempting, and wrong --- Zumdahl says so explicitly}

\examsection{II}{Free Response}{2 questions \textbullet\ 20 minutes}{\calculatorok}

% ---- FRQ 1 -----------------------------------------------------------------
\frq{short}{4}{Zumdahl \S11.1 \textbullet\ solution composition}

A solution is prepared by dissolving \SI{18.0}{\gram} of glucose
(\ce{C6H12O6}, $M = \SI{180.16}{\gram\per\mole}$) in \SI{182.0}{\gram} of
water. The resulting solution has a volume of \SI{190.}{\milli\liter}.

\begin{parts}
  \item Calculate the molality of the solution. \workspace[2cm]
        \keyonly{\begin{apcanswer}$n = 18.0/180.16 = \SI{0.0999}{\mole}$;
        $m = 0.0999/0.1820 =
        \SI{0.549}{\mole\per\kilo\gram}$\end{apcanswer}}
  \item Calculate the molarity of the solution. \workspace[1.8cm]
        \keyonly{\begin{apcanswer}$M = 0.0999/0.190 =
        \SI{0.526}{\Molar}$\end{apcanswer}}
  \item The solution is warmed to \SI{60}{\celsius}. State what happens to
        each of your two answers, with a reason.
        \answerlines[3]{The molality is unchanged --- both moles of solute
        and mass of solvent are independent of temperature. The molarity
        decreases --- the solution expands, so the same moles occupy a
        larger volume.}
\end{parts}

\begin{rubric}
  \rpoint{1}{(a) uses \SI{0.1820}{\kilo\gram} of SOLVENT --- using the
    solution mass earns 0}
  \rpoint{1}{(a) \SI{0.549}{\mole\per\kilo\gram}}
  \rpoint{1}{(b) \SI{0.526}{\Molar} using the solution volume}
  \rpoint{1}{(c) both effects correct WITH reasons}
\end{rubric}

% ---- FRQ 2 -----------------------------------------------------------------
\frq{short}{4}{Zumdahl \S11.2--11.3 \textbullet\ energetics and solubility}

Two salts are dissolved separately in water. Salt X causes the solution to
warm; salt Y causes it to cool.

\begin{parts}
  \item State the sign of $\Delta H_{\text{soln}}$ for each and justify from
        the observation. \answerlines[2]{X: negative --- the solution warmed,
        so the dissolving released energy. Y: positive --- the solution
        cooled, so the dissolving absorbed energy from it.}
  \item Using the three-step model, explain what must be true of the
        hydration energy for salt Y. \answerlines[3]{Its hydration energy
        (step 3) is smaller in magnitude than the combined endothermic cost
        of breaking the lattice (step 1) and expanding the solvent (step 2),
        so the net enthalpy change is positive.}
  \item A student claims salt Y cannot dissolve because dissolving it
        requires energy. Evaluate this claim. \answerlines[3]{The claim is
        wrong. Enthalpy is not the only factor --- mixing increases disorder,
        and that increase can drive dissolving even when energy must be
        absorbed. Instant cold packs work on exactly this principle.}
\end{parts}

\begin{rubric}
  \rpoint{1}{(a) both signs correct, each tied to the observation}
  \rpoint{2}{(b) 1 pt names step 3 as smaller than steps 1+2 combined;
    1 pt states that this makes $\Delta H_{\text{soln}}$ positive}
  \rpoint{1}{(c) rejects the claim AND cites disorder/entropy --- ``it just
    does'' earns 0}
\end{rubric}

\examtotal

\end{document}
